\documentclass[11pt]{article}
\usepackage[a4paper,margin=24mm]{geometry}
\usepackage{amsmath,amssymb,amsthm,booktabs,longtable,hyperref}
\newtheorem{theorem}{Theorem}\newtheorem{lemma}{Lemma}
\newcommand{\PG}{\operatorname{PG}}\newcommand{\PSL}{\operatorname{PSL}}
\title{Two literal projective-plane constructions for $\PSL_3(4)$}
\author{A bounded recovery of constructions 5.1 and 5.2 in Gottschalk--Leemans}
\date{}
\begin{document}\maketitle
\noindent
The reference edition is the complete 22-page Halle preprint \cite{GL}, not
an uncollated later journal edition. We recover the two incidence
constructions on its physical pages 14--16, Lemmas 4.1--4.3 and Table 1.
The finite certificate also treats the specified truncations below. This is
an existence and incidence result, not an exhaustion of higher-rank geometries.
Automorphism and correlation groups, conjugacy-fusion counts, doubling and
the other constructions are not asserted here.

\paragraph{The actual finite model.}
Put $F=\mathbb F_2[w]/(w^2+w+1)$ and label its elements $0,1,2=w,3=w+1$.
List the 21 projective points by lexicographically sorting the nonzero
column vectors whose first nonzero coordinate is 1. Lines are the kernels
of the 21 normalized nonzero linear forms. A hyperoval is a six-subset with
no three collinear points. Let $G=\PSL_3(4)$ in its point action.

We use the two actual degree-21 generators retained in the independently
reviewed rank-two component. No equality of independently ordered labels
is presumed. A five-point orbit of its order-960 line stabilizer gives a
line; the orbit of that set gives all 21 lines. The certificate verifies the
coordinate-to-retained label map
\[
(6,2,10,14,18,1,7,8,9,3,11,16,21,4,19,12,17,5,15,20,13).
\]
It checks every line under this bijection. In these coordinates the actual
two retained generators are represented, respectively, by
\[
 \begin{pmatrix}1&0&0\\0&w&0\\0&0&w+1\end{pmatrix},\qquad
 \begin{pmatrix}0&0&1\\1&0&0\\0&1&1\end{pmatrix}.                 \tag{1}
\]
Both determinants are 1; all 21 images are checked, not merely a projective
frame. Closure of the two literal permutations has 20,160 elements. This is
also $|\mathrm{SL}_3(4)|/3=4^3(4^3-1)(4^2-1)/3$, so the realization is
the full named projective special linear group. The retained named-model
identification and finite-field conventions are explicit inputs; the literal
matrix and permutation checks additionally verify their use here.

Testing all $\binom{21}{6}=54,264$ subsets gives 168 hyperovals, in three
$G$-orbits of 56. Fix distinct orbits $\mathcal H_3,\mathcal H_4$.
The five types are lines, points, $\mathcal H_3$, $\mathcal H_4$, and
unordered point pairs. The type sizes are $(21,21,56,56,210)$.
For $h\in\mathcal H_3$, $k\in\mathcal H_4$ and $Q$ a point pair, define
the ten relations by
\[
\begin{array}{c|cccc}
 &p&h&k&Q\\\hline
\ell&p\notin\ell&\ell\cap h=\varnothing&
 \ell\cap k\ne\varnothing&\ell\cap Q=\varnothing\\
p&&p\in h&p\in k&p\in Q\\
h&&&|h\cap k|=3&Q\subset h\\
k&&&&Q\subset k .
\end{array}                                                     \tag{2}
\]
This is $\Gamma_{5.1}$. For $\Gamma_{5.2}$ replace only
$\ell\cap k\ne\varnothing$ by $\ell\cap k=\varnothing$.
These are the literal rules in the source; in particular line--point
incidence is nonincidence in the underlying projective plane.

\begin{theorem}
Each of these two incidence systems is a firm rank-five geometry with
20,160 chambers. The actual $G$ acts regularly on chambers and transitively
on flags of every type. Each geometry is residually connected, residually
weakly primitive, and satisfies $\mathrm{IP}_2$. Its rank-one residue
double-transitivity flags, in type order, are
$(\mathrm{false},\mathrm{true},\mathrm{true},\mathrm{true},\mathrm{true})$.
For $\Gamma_{5.1}$ all twenty one-type
extension numbers in source Table 1 are those displayed in Table~\ref{tab:ext}.
The choice of the ordered pair of distinct hyperoval orbits does not change
either geometry up to a type-preserving isomorphism normalizing the
specified $G$-action.
\end{theorem}

\paragraph{Exact finite verification and its completeness.}
The accompanying program uses only integer masks, permutations and finite
field arithmetic after exporting the fixed input generators and line orbit.
For a prescribed ordered type set $I$, it recursively chooses an object of
each type, intersecting the remaining allowable sets with its adjacency
row. Induction on $|I|$ proves that the returned tuples are precisely all
flags of those types, without repetition. This procedure is applied to
every subset of the five types, including the empty subset.

Let $\mathcal C$ be the resulting complete set of full flags and choose
$C\in\mathcal C$. Applying every one of the 20,160 literal elements of
$G$ to $C$ gives 20,160 distinct tuples, exactly $\mathcal C$. All generator
actions preserve every object set and all ten incidence relations. Thus
this is a genuine regular chamber action; its kernel is trivial, since
an element in the kernel fixes $C$. Independently, for every type subset
$I$ the complete flag list is compared with the projection of $\mathcal C$
onto $I$. Equality proves that every partial flag extends to a chamber;
chamber transitivity alone is not substituted for this check. It then also
proves flag transitivity.

For each $F\subseteq C$, the residue objects of type $i$ are computed as
the intersection of the adjacency rows of the objects of $F$. Its action
uses precisely those literal elements of $G$ fixing every object of $F$.
No permutation-group quotient or replacement action is used. The incidence
graph of each residue of rank at least two is connected by direct traversal.
Each panel has at least two extensions, proving firmness under the standard
panel convention; a chamber itself is not required to have two extensions.
Because every flag is conjugate to the corresponding subflag of $C$, these
are checks of all residues, including those based at flags not chosen in
the calculation.

\begin{lemma}
Let $K$ act transitively on a finite set $X$, with $x\in X$. For one
representative $y$ of each $K_x$-orbit on $X\setminus\{x\}$, form the
undirected graph with edges $\{x^g,y^g\}$ for all $g\in K$.
The action is primitive if and only if all these graphs are connected.
\end{lemma}
\begin{proof}
Connected components of any such graph are a $K$-invariant equivalence
relation, so a disconnected graph contradicts primitivity. Conversely, if
a nontrivial block contains $x$ and $y\ne x$, all edges in that pair orbit
lie within blocks; its graph cannot be connected. Choosing another point
in the same $K_x$-orbit gives the same graph.
\end{proof}

The certificate applies this lemma to the actual actions of each $G_F$
on each remaining type. Every residue has at least one primitive type
action, precisely the weak primitivity required by RWPRI. Double
transitivity is recorded separately by the orbits of an actual point
stabilizer; it is never a filter. For a rank-two residue, the program
checks that the graph is complete bipartite or that each pair on one side
has at most one common neighbour. The latter is symmetric (a pair of
common neighbours would give a quadrangle) and proves the partial-linear
condition $\mathrm{IP}_2$. These checks use every pair of local objects.

\begin{table}[t]\centering
\begin{tabular}{ccc@{\qquad}ccc}\toprule
fixed types&added type&number&fixed types&added type&number\\\midrule
123&4&5&123&5&5\\124&3&3&124&5&3\\125&3&2&125&4&2\\
134&2&3&134&5&3\\135&2&2&135&4&2\\145&2&2&145&3&2\\
234&1&3&234&5&2\\235&1&6&235&4&4\\245&1&6&245&3&4\\
345&1&3&345&2&2\\\bottomrule
\end{tabular}
\caption{Exact source Table 1 extension numbers for $\Gamma_{5.1}$.}
\label{tab:ext}\end{table}

The finite output retains all 32 flag counts, all actual stabilizer orders,
all 31 positive-rank residue records, every local primitivity and
double-transitivity flag, and every rank-two test for each model. In both
models the panel orders in type order are $(2,1,1,1,1)$; equivalently the
five panel sizes are $(3,2,2,2,2)$. Their chamber stabilizers are trivial.
The first panel action is the regular cyclic group of order 3 and is not
doubly transitive; the other four panel actions have degree 2.
Table~\ref{tab:ext} is computed directly at the corresponding base triple
flags; the already verified flag action transports it to every triple flag.

\paragraph{The six choices of hyperoval orbits.}
The coordinate maps $\operatorname{diag}(w,1,1)$ and
$(x,y,z)\mapsto(x^2,y^2,z^2)$ preserve the plane and normalize the actual
group: conjugates of both retained generators are checked in the complete
literal $G$. Their permutations of the three hyperoval orbits are a
3-cycle and a transposition. They therefore act transitively on the six
ordered pairs of distinct orbits. For each of these six pairs, for both
models, and for each of the two maps, the certificate stores the resulting
bijections on all five object types and checks every incidence and
nonincidence. These 24 complete maps pay the choice-independence assertion.
They are not a computation of the full automorphism group.

\paragraph{The specified truncations.}
The literal restrictions in Table~\ref{tab:trunc} have the indicated source
profiles. Here the type numbers always refer to (2), $B$ is the chamber
stabilizer in the specified $G$, and panel orders are listed in increasing
retained-type order. For each row the program again enumerates every partial
flag and chamber, proves flag transitivity, and verifies every base residue,
firmness, RWPRI and $\mathrm{IP}_2$. It checks that projection of all
rank-five chambers is exactly the truncated chamber set.

\begin{table}[h]\centering
\begin{tabular}{ccclc}\toprule
source label&parent&retained types&panel orders&$|B|$\\\midrule
4.1&5.1, 5.2&2345&$(1,3,3,1)$&3\\
4.10&5.1&1345&$(2,1,1,2)$&2\\
3.1&5.1, 5.2&235&$(1,3,4)$&12\\
3.6&5.1&345&$(3,3,2)$&6\\
2.2&5.1, 5.2&35&$(3,14)$&24\\\bottomrule
\end{tabular}
\caption{Eight checked literal truncation routes.}\label{tab:trunc}
\end{table}

The object maps are the identity on each retained type; incidence is
literally restricted, rather than inferred from the numerical profile.
The two 4.1 restrictions and the two 3.1 restrictions are therefore
identical incidence models: they omit the only relation changed between
5.1 and 5.2. In particular the former also restricts to the latter and to
3.6. The 4.10 restriction likewise restricts to 3.6. These maps establish
the advertised construction routes in this fixed edition. They do not
identify an arbitrary geometry from its diagram alone. No further
independent realization from the cited earlier papers is being asserted
isomorphic without a supplied map.

The rank-two 2.2 restriction has local degrees $(15,4)$ with parabolic
orders $(360,96)$ and intersection order 24. Its actual local
double-transitivity flags are $(\mathrm{false},\mathrm{true})$, consistent
with the corrected accepted rank-two component: $14\nmid24$. Thus the
local-2T distinction is preserved through the higher-rank construction.

\paragraph{Scope and attribution.}
The constructions and their intended geometric properties are due to
Gottschalk and Leemans and to the predecessors explicitly cited in their
source. The new material here is a transparent finite certificate of those
specified constructions, together with the exact action and source joins.
The accepted rank-two model is acquired unchanged. Complete higher-rank
classification, other incidence constructions, doubling, full group names
and kernels beyond the orders stated above, full automorphism and
correlation groups, and fusion multiplicities remain separate obligations.
No novelty, whole-source, question, Lean or measurement credit is claimed.

\begin{thebibliography}{9}
\bibitem{GL} H.~Gottschalk and D.~Leemans, \emph{Geometries for the group
$\mathrm{PSL}(3,4)$}, Halle Reports on Algebra and Geometry 99--30,
fixed 22-page preprint; physical pp.~8--16 for the source routes used here.
\bibitem{R2} \emph{Primitive residues for $\mathrm{PSL}_3(4)$}, accepted
L100 rank-two component, edition 2, with distinct review and root adoption;
exact immutable manifests are listed in the certificate dependencies.
\end{thebibliography}
\end{document}
